
Metadata repositories for machine learning datasets---HuggingFace Hub, OpenML, and UCI Machine Learning Repository---employ heterogeneous documentation schemas that prevent systematic cross-repository analysis. Researchers conducting systematic reviews or meta-analyses across repositories must manually map fields to identify which metadata addresses responsible AI concerns versus general dataset characteristics. This process does not scale to the thousands of datasets now available.

The Datasheets for Datasets framework \citep{Gebru2018} introduced lifecycle-stage categories for dataset documentation, including motivation, composition, collection process, and responsible AI considerations. Roman et al. (2023) operationalized these categories in a two-tier structure (General Information versus Responsible AI) within HuggingFace's Open Datasheets interface. However, this approach requires explicit template adoption within a single repository and does not address the inverse problem: discovering lifecycle structure in heterogeneous metadata that lacks standardized templates.

This study investigated whether semantic embeddings of metadata fields encode lifecycle categories as geometric structure accessible to unsupervised clustering. If successful, this would enable automated cross-repository mapping without manual annotation or template enforcement. We tested two specific hypotheses: (1) lifecycle categories exhibit operational stability across repositories and linear separability in embedding space (signal existence), and (2) unsupervised clustering recovers lifecycle structure with normalized mutual information exceeding 0.60 and outperforming baseline methods by at least 0.15 NMI (unsupervised recovery).

The results revealed a striking asymmetry. Linear probes trained on 60 samples achieved 79.6\% accuracy across repositories, confirming that embeddings encode lifecycle information. However, K-means clustering achieved NMI of 0.023, barely exceeding random permutation (0.010) and failing to reach the 0.60 threshold. Lexical keyword matching for terms like ''privacy,'' ''ethics,'' and ''license'' detected zero Responsible AI fields, indicating extreme terminology variation across repositories.

We attribute clustering failure to severe class imbalance (8.3\% Responsible AI fields) compounded by repository heterogeneity. Repository-stratified analysis showed UCI metadata achieving NMI of 0.394 (still below threshold but 30 times higher than HuggingFace's 0.013), correlating with UCI's less extreme class imbalance (14\% Responsible AI versus HuggingFace's 5.3\%). These findings indicate that lifecycle separability exists as a supervised signal requiring task-specific amplification rather than as natural cluster structure.

The practical implication is that cross-repository lifecycle detection requires supervised or semi-supervised methods---few-shot learning, active learning, or label propagation---rather than unsupervised discovery. This study provides preliminary evidence that 60 training samples enable 76\% generalization to a held-out repository, suggesting feasibility for label-efficient deployment, though production-scale validation remains necessary.

